\chapter{Preliminary Results}
\label{chap:results}

Each section states the pre-registered question, the result and its status. All numbers are on held-out questions; confidence intervals are 95\% paired bootstrap intervals and $p$ values come from exact McNemar tests, unless stated otherwise.

\section{A small selector saves most of the tokens (HotpotQA)}
\label{sec:res_selector}

A1, the fine-tuned 22.7M cross-encoder, was compared with sending all ten paragraphs ($k=10$) on the 400 HotpotQA test questions. Selection rules were tuned on the development set. The reader was \texttt{gpt-5.4-mini} with the prompt of \citeonline{lavado2027noisegate}. Table~\ref{tab:hotpot} shows that all A1 rules reach answer F1 statistically indistinguishable from $k=10$, with 68--77\% fewer prompt tokens. The gold-paragraphs oracle is better than $k=10$ by 0.040 F1 (CI $[0.017, 0.064]$), so distractors actively hurt the reader.

\begin{table}[htbp]
\centering
\caption{HotpotQA test (400 questions): answer quality and prompt tokens. $\Delta$F1 is paired against $k=10$.}
\label{tab:hotpot}
\small
\begin{tabular}{lrrrl}
\toprule
Context & F1 & EM & Tokens & $\Delta$F1 [95\% CI] \\
\midrule
All 10 paragraphs ($k=10$) & 0.782 & 0.620 & 1,448 & --- \\
A1, threshold for recall 0.95 & 0.796 & 0.630 & 467 ($-68\%$) & $+0.014$ $[-0.013, +0.043]$ \\
A1, threshold for recall 0.90 & 0.778 & 0.623 & 397 ($-73\%$) & $-0.004$ $[-0.033, +0.026]$ \\
A1, top-2 & 0.780 & 0.628 & 330 ($-77\%$) & $-0.001$ $[-0.033, +0.030]$ \\
Gold paragraphs (oracle) & 0.822 & 0.670 & 329 & $+0.040$ $[+0.017, +0.064]$ \\
\bottomrule
\end{tabular}
\end{table}

\section{Adaptive stopping saves further tokens}
\label{sec:res_adaptive}

On HotpotQA, choosing between two and four paragraphs per question with a learned stopping rule saved 57 and 71 tokens per question against A1's threshold at recall targets 0.90 and 0.95 (CIs $[-71, -44]$ and $[-91, -53]$), with recall differences not significant. On MuSiQue, adaptive stopping on top of a ranker fine-tuned in domain reached recall 0.95 with 1,554 tokens and 20 passes per question, against 2,535 tokens for sending all paragraphs. The largest-gap rule Adaptive-$k$ \cite{taguchi2025adaptivek}, applied to the same cross-encoder scores, selected 3.7 paragraphs on average with recall 0.50; with a buffer of five extra paragraphs it spent the same tokens as the learned rule but recovered fewer complete chains (0.815 vs 0.868; exploratory). This agrees with its failure on cosine scores in \citeonline{lavado2027noisegate}.

\section{On MuSiQue the ranker is the bottleneck}
\label{sec:res_bottleneck}

With A1 (trained on HotpotQA), the top ten paragraphs contained all supporting paragraphs for only 68\% of MuSiQue questions, so no stopping rule limited to $k\le10$ could reach the planned recall targets; the pre-registered comparison was declared not testable. Fine-tuning A1 on 1,000 MuSiQue questions (M1) raised full-support recall@5 from 0.355 to 0.682, and reduced the tokens needed for recall 0.90 from 1,957 to 1,533. The iterative model M1-iter reached 0.750 (McNemar $p=0.0067$ against M1), with the largest gain on four-hop questions (0.58 vs 0.36). That comparison, however, was confounded by the amount of training data (Section~\ref{sec:res_matched}).

\section{Chain-following is a rank-2 change}
\label{sec:res_rank}

The weight difference $\Delta = \text{M1-iter} - \text{M1}$ was grafted onto M1 at ranks 1 to 16 (Equation~\ref{eq:recovery}), and the model was run iteratively. Without the graft, M1 run iteratively reached recall@5 of 0.45 on \texttt{dev300}, against 0.74 for M1-iter. Recovery was 74\% at rank 1 and 95\% at rank 2, and stayed flat at higher ranks (97\% at ranks 4 and 8). A random rank-4 change of the same norm recovered 9\%. On the test set, the rank-2 graft recovered 84.7\% of the gain (CI $[73.6\%, 96.2\%]$); it remained below the full model (14 vs 31 discordant questions, $p=0.016$). Grafting single layers showed that layers 3 and 4 contribute most (43\% and 40\%), but layers 0--2 also contribute about 30\% each, a flatter profile than for the HotpotQA skill (A1 $-$ A0), where a rank-4 graft recovered 80\% of a token-based measure of the gain and the middle layers dominated. Ranks 2 to 8 differ by about two questions on \texttt{dev300}; the result is a plateau from rank 2, not a strict ordering.

\section{Pair-matched control}
\label{sec:res_matched}

Figure~\ref{fig:matched} illustrates the design. With identical training pairs, M1-iter run iteratively beat M1-ctrl run one-shot by 0.040 recall@5 on the test set (0.7475 vs 0.7075, CI $[-0.005, 0.085]$, $p=0.11$): the pre-registered primary test failed. The difference was positive for all three seeds ($+0.040$, $+0.028$, $+0.033$), and the minimum detectable effect was about 0.07, so the result indicates a small effect rather than none. Training on more pairs alone accounted for about 0.025 of the original gain (M1-ctrl vs M1, $p=0.013$, descriptive).

\begin{figure}[htbp]
\centering
% Pair-matched control design (TikZ). Requires \usepackage{tikz} and \usetikzlibrary{positioning,arrows.meta}.
\begin{tikzpicture}[
  box/.style={draw, rounded corners, align=center, minimum width=3.4cm, minimum height=1.0cm, font=\small},
  arr/.style={-{Latex[length=2mm]}, thick},
  node distance=0.7cm and 1.4cm]
  \node[box, fill=gray!10] (a1) {A1\\(HotpotQA selector)};
  \node[box, below left=of a1] (iter) {M1-iter\\query: question + selected};
  \node[box, below right=of a1] (ctrl) {M1-ctrl\\query: question only};
  \node[box, below=of iter, fill=blue!8] (ei) {iterative inference\\74 passes};
  \node[box, below=of ctrl, fill=blue!8] (ec) {one-shot inference\\20 passes};
  \draw[arr] (a1) -- node[left, font=\footnotesize, align=right] {same 18,101 pairs,\\same order} (iter);
  \draw[arr] (a1) -- (ctrl);
  \draw[arr] (iter) -- (ei);
  \draw[arr] (ctrl) -- (ec);
  \draw[<->, thick, dashed] (iter) -- node[above, font=\footnotesize] {$\Delta$ (rank-2 graft)} (ctrl);
  \draw[<->, thick] (ei) -- node[above, font=\footnotesize] {primary test} (ec);
\end{tikzpicture}

\caption{Pair-matched control. Both models start from A1 and see the same paragraphs, labels and order; only the query differs. Each is evaluated in the inference mode it was trained for, and their weight difference $\Delta$ is analysed by grafting.}
\label{fig:matched}
\end{figure}

The pre-registered secondary test asked whether the pure conditioning change, $\Delta = \text{M1-iter} - \text{M1-ctrl}$, is also low-rank. Grafted onto M1-ctrl at rank 2, it recovered 86.8\% of the effect on the test set (CI $[77.7\%, 95.8\%]$), while a random rank-2 change of the same norm recovered 7\%. Without this change, M1-ctrl collapses when run iteratively (recall@5 of 0.41), so $\Delta$ encodes the ability to \emph{read} a conditioned query.

\section{The effect is concentrated on long chains}
\label{sec:res_long}

An exploratory reading of the matched control suggested that the gain was concentrated on four-hop questions. This was tested on 566 fresh four-hop questions (166 unused development questions and 400 unused training questions). M1-iter (iterative) beat M1-ctrl (one-shot) by 0.055 recall@5 (0.500 vs 0.445, CI $[0.007, 0.102]$, $p=0.035$), positive in both sources ($+0.139$ and $+0.020$) and for all three seeds. On all 1,017 fresh development questions (exploratory), the gain was about zero on two-hop questions, $+0.078$ on three-hop ($p=0.010$) and $+0.139$ on four-hop questions ($p=0.006$).

A cascade that runs the one-shot ranker for every question and the iterative one only for questions predicted to be four-hop did not help: it routed 29\% of questions, reached 0.702 recall@5 with 41.5 passes, against 0.693 for one-shot (20 passes) and 0.745 for always-iterative (74 passes), and failed both pre-registered tests. The router predicted hop count reasonably (four-hop recall 0.76, precision 0.43), but not the questions where iteration actually helps.

\section{Answer quality}
\label{sec:res_f1}

\todo[inline]{exp15 (running): answer F1 of gpt-5.4-mini with the top-5 sets of M1-ctrl and M1-iter on the fresh development questions; pre-registered primary on questions with three or more hops, order-controlled.}

\section{Summary}

Table~\ref{tab:claims} lists the claims and their status. The practical result is that a small in-domain selector with adaptive stopping is the cost-efficient choice. The research result is that the ability to follow an evidence chain is a separable rank-2 change in such a selector, with a measurable benefit on long chains only, and that simple routing does not yet capture that benefit.

\begin{table}[htbp]
\centering
\caption{Claims and their status (C = confirmed by a pre-registered test; N = null; E = exploratory).}
\label{tab:claims}
\small
\begin{tabularx}{\textwidth}{>{\raggedright\arraybackslash}Xll}
\toprule
Claim & Status & Section \\
\midrule
Small selector: 68--77\% fewer tokens at equal F1 (HotpotQA) & C & \ref{sec:res_selector} \\
Adaptive stopping saves tokens at a recall target & C & \ref{sec:res_adaptive} \\
Iteration beats one-shot at matched training data & N ($+0.040$, $p=0.11$) & \ref{sec:res_matched} \\
The conditioning change is rank 2 & C (86.8\%) & \ref{sec:res_rank}, \ref{sec:res_matched} \\
Iteration helps four-hop chains & C ($+0.055$) & \ref{sec:res_long} \\
Iteration helps three-hop chains, not two-hop & E & \ref{sec:res_long} \\
A hop-count router gives a useful cascade & N & \ref{sec:res_long} \\
\bottomrule
\end{tabularx}
\end{table}
